\documentclass[11pt,a4paper]{article}
\usepackage[T2A]{fontenc}
\usepackage[utf8]{inputenc}
\usepackage[russian,english]{babel}
\usepackage{amsmath,amssymb,amsthm,mathtools}
\usepackage{setspace}
\usepackage{enumitem}
\usepackage[hidelinks]{hyperref}
\onehalfspacing
\renewcommand{\leq}{\leqslant}
\renewcommand{\le}{\leqslant}
\begin{document}
\begin{center}
{\huge\bfseries\foreignlanguage{russian}{Подстановка $x=2\cos\theta$ и\\ многочлены Чебышёва}\par}\vspace{1.2em}
{\Large\foreignlanguage{russian}{Домашнее задание}\par}\vspace{0.8em}
{\foreignlanguage{russian}{1 сентября 2026 года}\par}
\end{center}\vspace{0.5em}

\noindent\textbf{\foreignlanguage{russian}{Теоретический минимум.}}\quad
$T_n(\cos\theta)=\cos n\theta$, $U_{n-1}(\cos\theta)=\sin n\theta/\sin\theta$, the recurrence
$T_{n+1}=2xT_n-T_{n-1}$ and the hyperbolic branch $T_n(\cosh t)=\cosh nt$; Chebyshev's extremal
theorem; the inequalities of Markov and Bernstein; orthogonality with weight $(1-x^2)^{-1/2}$; the
semicircle substitution $x=2\cos\theta$ in three-term recurrences.

\medskip
\noindent\rule{\textwidth}{0.4pt}
\medskip

\begin{enumerate}[itemsep=0.6em]
\item Prove that $T_m\circ T_n=T_{mn}$ and that these are, up to affine conjugacy, the only
polynomial families of degree $>1$ closed under composition together with the powers $x^n$
(\foreignlanguage{russian}{теорема Ритта, достаточно сформулировать и проверить оба примера}).

\item Prove the orthogonality relations $\int_{-1}^1\frac{T_mT_n}{\sqrt{1-x^2}}\,dx=0$ for $m\ne n$,
and expand $x^n$ in the basis $T_0,\dots,T_n$.

\item Prove that the roots of $U_{n-1}$ are $\cos\frac{k\pi}{n}$ and deduce that the eigenvalues of
the tridiagonal matrix with $0$ on the diagonal and $1$ on the neighbouring diagonals are
$2\cos\frac{k\pi}{n+1}$; compare with the discrete Wirtinger inequality from the first sheet.

\item Prove that $\min_p\max_{[a,b]}|p|=2\bigl(\frac{b-a}{4}\bigr)^n$ over monic $p$ of degree $n$,
and find the polynomial attaining it.

\item (Proposed problems, \foreignlanguage{russian}{№}235) Prove that $\cos\frac{\pi}{7}$ has minimal
polynomial $x^3-\frac12x^2-\frac12x+\frac18$ over $\mathbb Q$, that $\cos\frac{2\pi}{7}$ has minimal
polynomial $x^3+\frac12x^2-\frac12x-\frac18$, and that $\cos\frac{3\pi}{7}$ has the same minimal
polynomial as $\cos\frac{\pi}{7}$.

\item Prove that $2\cos\frac{2\pi}{n}$ is a rational number only for $n\in\{1,2,3,4,6\}$, using the
recurrence $T_{n+1}=2xT_n-T_{n-1}$ and the rational root theorem.

\item Let $p$ be a real polynomial of degree $n$ with $|p(x)|\le1$ on $[-1,1]$. Prove that
$|p(x)|\le|T_n(x)|$ for every $|x|>1$, and that the leading coefficient of $p$ is at most $2^{n-1}$
in absolute value.

\item Prove that $\sum_{k=0}^nT_k(x)t^k=\frac{1-tx}{1-2tx+t^2}+O(t^{n+1})$, i.e.\ find the generating
functions of $T_n$ and $U_n$, and deduce the addition formula $T_{m+n}+T_{m-n}=2T_mT_n$.

\item (Solve in radicals) Solve the equation $x^3-3x=c$ for every real $c$ by the substitution
$x=2\cos\theta$ when $|c|\le2$ and $x=2\cosh t$ otherwise; explain what the casus irreducibilis of
Cardano's formula looks like in this language.

\item (IMC 2026) Prove that there exists a constant $C>0$ such that for every pair $A$, $B$ of
positive integers there is a real polynomial $p(x)$ with
\[
p(0)^2>\sum_{i=1}^{A}p(-i)^2+\sum_{i=1}^{B}p(i)^2\qquad\text{and}\qquad\deg p<C\sqrt{AB}.
\]
\end{enumerate}
\end{document}
